\documentclass[11pt]{article}
\usepackage[margin=1in]{geometry}
\usepackage{amsmath,amssymb,amsthm,enumitem,hyperref}
\hypersetup{colorlinks=true,linkcolor=blue,urlcolor=blue,citecolor=blue}
\newtheorem{theorem}{Theorem}
\newtheorem{lemma}{Lemma}
\newcommand{\Z}{\mathbb Z}
\newcommand{\R}{\mathbb R}
\newcommand{\one}{\mathbf1}
\title{Ordinal-sum factorization for polar preorder polytopes\\\large MathDB 374964}
\author{Principia Math In Progress}
\date{September 18, 2026}
\begin{document}
\maketitle
\begin{abstract}
We prove that the $h^*$-polynomial of the polar reflexive preorder
polytope factors under ordinal sums, as conjectured by Athanasiadis
and Chapoton. An explicit piecewise linear bijection of integer
lattices makes the corresponding Minkowski functionals additive.
Their lattice generating functions, and hence their $h^*$-polynomials,
therefore multiply.
\end{abstract}

\section*{Task 1. Full solution}
For a preorder $\tau$ on a nonempty finite set $E$, let
\begin{equation}\label{eq:K}
 K_\tau=\operatorname{conv}\bigl(\{-e_i:i\in E\}
       \cup\{\one_I:\varnothing\ne I\text{ an order ideal of }\tau\}\bigr).
\end{equation}
Here $e_i$ is a standard basis vector and $\one_I$ is the
indicator vector of $I$. By \cite[Section 6]{AC}, this is precisely
the polar polytope $\mathcal R_\tau^\vee$ in the question.
An ordinal sum $\tau_1\oplus\tau_2$ on disjoint sets $E_1,E_2$
retains each internal preorder and declares every element of $E_1$
smaller than every element of $E_2$.

\begin{theorem}\label{thm:main}
For arbitrary finite preorders $\tau_1,\tau_2$,
\[
 \boxed{h^*(K_{\tau_1\oplus\tau_2},t)
       =h^*(K_{\tau_1},t)h^*(K_{\tau_2},t).}
\]
Thus Conjecture 6.3 of \cite{AC} holds.
\end{theorem}

\subsection*{An explicit Minkowski functional}
The order polytope in the decreasing convention is
\[
 O_\tau=\{u\in[0,1]^E:u_i\ge u_j\text{ whenever }i\le_\tau j\}.
\]
It equals the convex hull of the indicators of all order ideals,
including the empty ideal. For completeness, for any $u\in O_\tau$
each nonempty superlevel set $\{i:u_i\ge a\}$ is an ideal.
Writing $u$ as the sum of its successive level differences times
these indicators, and adding the remaining coefficient at $0$,
gives the required convex combination. The reverse inclusion is immediate.

The origin is in the interior of $K_\tau$: it is a convex
combination with strictly positive coefficients of the affinely
independent vertices $-e_i$ and $\one_E$. Consequently $K_\tau$
is full dimensional, and
\[
 K_\tau=\operatorname{conv}(O_\tau\cup(-\Delta_E)),
 \qquad \Delta_E=\operatorname{conv}(0,e_i:i\in E).
\]
For $z\in\R^E$, define
\[
 M(z)=\max(\{0\}\cup\{z_i:i\in E\}),\qquad
 U_i(z)=\max(\{0\}\cup\{z_j:i\le_\tau j\}).
\]
Let $g_\tau(z)=\inf\{k\ge0:z\in kK_\tau\}$ be its Minkowski
functional, also called its gauge.

\begin{lemma}\label{lem:gauge}
For every $z\in\R^E$,
\begin{equation}\label{eq:gauge}
 g_\tau(z)=M(z)+\sum_{i\in E}U_i(z)-\sum_{i\in E}z_i.
\end{equation}
In particular, $g_\tau$ is integer valued on $\Z^E$.
\end{lemma}
\begin{proof}
The convex-hull description says that $z\in kK_\tau$ exactly
when $z=u-v$, where $u$ is nonnegative and weakly decreasing
along $\tau$, $v$ is nonnegative, and
\[
 \max_i u_i+\sum_i v_i\le k.
\]
Indeed, $u$ belongs to $(\max_i u_i)O_\tau$ and $-v$ belongs
to $(\sum_i v_i)(-\Delta_E)$; this proves one direction.
Conversely a convex combination from the two polytopes gives
such a decomposition with these two quantities bounded by its
respective coefficients times $k$.

Equivalently, minimize $\max_i u_i+\sum_i(u_i-z_i)$ among
nonnegative decreasing vectors $u$ with $u\ge z$ coordinatewise.
The coordinatewise least such vector is $U(z)$. Reflexivity gives
$U_i(z)\ge z_i$, and transitivity gives $U_i(z)\ge U_j(z)$
when $i\le_\tau j$. Every other feasible $u$ satisfies
$u_i\ge u_j\ge z_j$ for all $j$ above $i$, so $u_i\ge U_i(z)$.
The objective is nondecreasing in every coordinate of $u$ and
therefore takes its minimum at $U(z)$. Finally
$\max_i U_i(z)=M(z)$, proving \eqref{eq:gauge}.
\end{proof}

\subsection*{The lattice bijection for an ordinal sum}
Put $\tau=\tau_1\oplus\tau_2$ and write $z=(x,y)$ according
to $E_1\sqcup E_2$. Set $c=M(y)$. Define
\begin{equation}\label{eq:phi}
 \Phi(x,y)=(x-c\one_{E_1},y).
\end{equation}
This is a bijection of $\Z^{E_1}\times\Z^{E_2}$, with inverse
$(x',y)\mapsto(x'+M(y)\one_{E_1},y)$.
We claim the stronger identity
\begin{equation}\label{eq:add}
 g_\tau(x,y)=g_{\tau_1}(x-c\one_{E_1})+g_{\tau_2}(y)
 \qquad(x,y\text{ real}).
\end{equation}
Write $U_i^1(x)$ and $U_j^2(y)$ for the maxima within the two
preorders. The maxima for the ordinal sum are
\[
 U_i(x,y)=\max(U_i^1(x),c)\quad(i\in E_1),\qquad
 U_j(x,y)=U_j^2(y)\quad(j\in E_2).
\]
Since $c\ge0$, with $a_+=\max(a,0)$ one also has
\[
 U_i^1(x-c\one)=(U_i^1(x)-c)_+,
 \qquad M(x-c\one)=(M(x)-c)_+.
\]
If $p=|E_1|$, formula \eqref{eq:gauge} now gives
\begin{align*}
 g_{\tau_1}(x-c\one)
 &= (M(x)-c)_++\sum_{i\in E_1}(U_i^1(x)-c)_+
                  -\sum_{i\in E_1}x_i+pc\\
 &=\max(M(x),c)-c+\sum_{i\in E_1}\max(U_i^1(x),c)
                  -\sum_{i\in E_1}x_i.
\end{align*}
Adding $g_{\tau_2}(y)=c+\sum_j U_j^2(y)-\sum_j y_j$
is exactly \eqref{eq:gauge} for $\tau$, proving \eqref{eq:add}.

\subsection*{Multiplication of the generating functions}
Define the formal series
\[
 \Theta_\tau(t)=\sum_{z\in\Z^E}t^{g_\tau(z)}.
\]
Its exponents are nonnegative integers. Its coefficients are finite
because $\{z:g_\tau(z)\le k\}=kK_\tau$ is bounded.
The bijection \eqref{eq:phi} and identity \eqref{eq:add} imply
\[
 \Theta_{\tau_1\oplus\tau_2}(t)
 =\Theta_{\tau_1}(t)\Theta_{\tau_2}(t).
\]
For $n=|E|$, summing over all dilates containing each lattice point
gives
\[
 \sum_{k\ge0}\#(kK_\tau\cap\Z^E)t^k
 =\frac{\Theta_\tau(t)}{1-t}
 =\frac{h^*(K_\tau,t)}{(1-t)^{n+1}}.
\]
Thus $\Theta_\tau(t)=h^*(K_\tau,t)/(1-t)^n$.
Substitution into the multiplicative identity proves
Theorem~\ref{thm:main}.\hfill$\square$

The proof allows equivalent distinct elements in either preorder.
It also extends to an empty factor by assigning it the one-point
polytope and polynomial $1$. The map $\Phi$ is linear with an
integer matrix of determinant $1$ on each region where $M(y)$
is a fixed coordinate or zero. The argument requires this
piecewise linear lattice bijection, not a global linear equivalence.

\subsection*{Verification and attribution}
The polar vertex description and the conjecture are due to
Athanasiadis and Chapoton \cite{AC}. The explicit gauge formula
and lattice bijection supply the proof. The verifier compares the
gauge with independently enumerated rational facets, then checks
the bijection and full $h^*$-factorization for every ordinal pair
of total size at most four. Larger examples include equivalent
elements. All four exact verification groups pass.

All 26 pages of the current primary were read. Bounded checks
found no earlier full target claim or identified journal final.
The related support-enumerator duality and reflexive-polytope
word formulas concern different statements, as recorded in the
literature audit. Review is by the authoring agent only, without
independent peer review or proof-assistant certification.

\section*{Tasks 2--5. Proposed follow-up directions}
\begin{enumerate}[label=\textbf{Task \arabic*.},start=2,leftmargin=*]
\item Use the piecewise linear bijection to transport lattice
triangulations between the ordinal-sum polytope and the free sum
of the two polar polytopes.
\item Extend the gauge formula to weighted ideal indicators and
determine which integral weights preserve ordinal-sum factorization.
\item Find analogous lattice bijections for other operations on
preorders, including disjoint unions and substitutions at a vertex.
\item Study when the factorization of $h^*(K_\tau,t)$ detects
an ordinal-sum decomposition of the underlying preorder.
\end{enumerate}
\begin{thebibliography}{9}
\raggedright
\bibitem{AC} C. A. Athanasiadis and F. Chapoton,
\emph{Polytopes and posets associated to preorders},
\href{https://arxiv.org/abs/2605.26916v1}{arXiv:2605.26916v1}, 2026.
\end{thebibliography}
\end{document}
